\Article{21}
In the preceding calculations we left the exponent $n$ undetermined. On June $24$—$28$, $1832$, we carried out two series of experiments, extending them to the greatest distances the room permitted. They will show very clearly what value nature requires. In the first series the second needle was placed on a straight line perpendicular to the magnetic meridian (following the first method of Article $19$), and in the second series the center of the needle was placed on the meridian itself. The table gives an overview of the results of these experiments. The distances $R$ are expressed in fractions of a meter, and the values of the angle $\frac14(u-u'+u''-u''')$ are denoted by $v$ for the first series and by $v'$ for the second.
\begin{center}
\begin{tabular}{c|rrr|rrr}
$R$&\multicolumn{3}{c|}{$v$}&\multicolumn{3}{c}{$v'$}\\\hline
$1^{\mathrm m}1$&&&&$1^0$&$57'$&$24''8$\\
$1{,}2$&&&&$1$&$29$&$40{,}5$\\
$1{,}3$&$2^0$&$13'$&$51''2$&$1$&$10$&$19{,}3$\\
$1{,}4$&$1$&$47$&$28{,}6$&$0$&$55$&$58{,}9$\\
$1{,}5$&$1$&$27$&$19{,}1$&$0$&$45$&$14{,}3$\\
$1{,}6$&$1$&$12$&$7{,}6$&$0$&$37$&$12{,}2$\\
$1{,}7$&$1$&$0$&$9{,}9$&$0$&$30$&$57{,}9$\\
$1{,}8$&$0$&$50$&$52{,}5$&$0$&$25$&$59{,}5$\\
$1{,}9$&$0$&$43$&$21{,}8$&$0$&$22$&$9{,}2$\\
$2{,}0$&$0$&$37$&$16{,}2$&$0$&$19$&$1{,}6$\\
$2{,}1$&$0$&$32$&$4{,}6$&$0$&$16$&$24{,}7$\\
$2{,}5$&$0$&$18$&$51{,}9$&$0$&$9$&$36{,}1$\\
$3{,}0$&$0$&$11$&$0{,}7$&$0$&$5$&$33{,}7$\\
$3{,}5$&$0$&$6$&$56{,}9$&$0$&$3$&$28{,}9$\\
$4{,}0$&$0$&$4$&$35{,}9$&$0$&$2$&$22{,}2$
\end{tabular}
\end{center}

Even a glance at these figures shows that for larger values the numbers in the second column are approximately twice those in the third, and the numbers in both columns are approximately inversely proportional to the cube of the distances. Thus no doubt can remain that the value $n=2$ is correct. To confirm this \SourceLeaf{114}law still further in the individual experiments, we applied the method of least squares to all the figures and obtained the following coefficients.
\[
\begin{aligned}
\operatorname{tang}v&=0{,}086870R^{-3}-0{,}002185R^{-5}\\
\operatorname{tang}v'&=0{,}043435R^{-3}+0{,}002449R^{-5}
\end{aligned}
\]

Here is a comparison of the values calculated from these formulas with the observed values.
\begin{center}
\textit{Calculated values.}\par
\nopagebreak[4]
\setlength{\tabcolsep}{4pt}
\begin{tabular}{c||rrr|r||rrr|r}
$R$&\multicolumn{3}{c|}{$v$}&Difference&\multicolumn{3}{c|}{$v'$}&Difference\\\hline
$1^{\mathrm m}1$&&&&&$1^0$&$57'$&$22''0$&$+2''8$\\
$1{,}2$&&&&&$1$&$29$&$46{,}5$&$-6{,}0$\\
$1{,}3$&$2^0$&$13'$&$50''4$&$+0''8$&$1$&$10$&$13{,}3$&$+6{,}0$\\
$1{,}4$&$1$&$47$&$24{,}1$&$+4{,}5$&$0$&$55$&$58{,}7$&$+0{,}2$\\
$1{,}5$&$1$&$27$&$28{,}7$&$-9{,}6$&$0$&$45$&$20{,}9$&$-6{,}6$\\
$1{,}6$&$1$&$12$&$10{,}9$&$-3{,}3$&$0$&$37$&$15{,}4$&$-3{,}2$\\
$1{,}7$&$1$&$0$&$14{,}9$&$-5.0$&$0$&$30$&$59{,}1$&$-1{,}2$\\
$1{,}8$&$0$&$50$&$48{,}3$&$+4{,}2$&$0$&$26$&$2{,}9$&$-3{,}4$\\
$1{,}9$&$0$&$43$&$14{,}0$&$+7{,}8$&$0$&$22$&$6{,}6$&$+2{,}6$\\
$2{,}0$&$0$&$37$&$5{,}6$&$+10{,}6$&$0$&$18$&$55{,}7$&$+5{,}9$\\
$2{,}1$&$0$&$32$&$3{,}7$&$+0{,}9$&$0$&$16$&$19{,}8$&$+4{,}9$\\
$2{,}5$&$0$&$19$&$2{,}1$&$-10{,}2$&$0$&$9$&$38{,}6$&$-2{,}5$\\
$3{,}0$&$0$&$11$&$1{,}8$&$-1{,}1$&$0$&$5$&$33{,}9$&$-0{,}2$\\
$3{,}5$&$0$&$6$&$57{,}1$&$-0{,}2$&$0$&$3$&$29{,}8$&$-1{,}0$\\
$4{,}0$&$0$&$4$&$39.6$&$-3{,}7$&$0$&$2$&$20{,}5$&$+1{,}7$
\end{tabular}
\end{center}

